% --- INICIO DE PLANTILLA PERSONALIZADA ---
\documentclass[oneside,11pt]{Latex/Classes/PhDthesisPSnPDF}

% --- Paquetes Originales ---
\usepackage{blindtext}
\usepackage{actuarialangle} 
\usepackage{tabularx}
\usepackage{float}
\usepackage{multirow, array}
\usepackage[usenames,dvipsnames,svgnames,table]{xcolor}
\usepackage{longtable}
\usepackage{latexsym,amsmath,amssymb,amsfonts}
\usepackage{enumitem}
\usepackage{xparse}
\usepackage{booktabs}
\usepackage[round, sort, numbers]{natbib}  
\usepackage{adjustbox}
\usepackage[utf8]{inputenc}
\usepackage{caption}
\usepackage{graphicx}
\usepackage{hyperref}
\usepackage{cleveref}

% Definición de colores
\definecolor{miblue}{RGB}{68, 114, 196}

% --- CAMBIO CRÍTICO 1: Usar \input para macros en el preámbulo ---
% \include da problemas antes del begin document. \input es seguro.
\input{Latex/Macros/MacroFile1}

% --- Definiciones de seguridad (por si la clase falla al cargar alguna variable) ---
\providecommand{\facultad}[1]{\def\lafacultad{#1}}
\providecommand{\escudofacultad}[1]{\def\elescudo{#1}}
\providecommand{\degree}[1]{\def\elgrado{#1}}
\providecommand{\director}[1]{\def\eldirector{#1}}
\providecommand{\lugar}[1]{\def\ellugar{#1}}
\providecommand{\degreedate}[1]{\def\lafecha{#1}}

% --- Configuración para bloques de código de R (Pandoc) ---

%%%%%%%%%%%%%%%%%%%%%%%%%%%%%%%%%%%%%%%%%%%%%%%%%%%%%%%%%%%%%%%%%%%%%%%%%%%%%%%%
%                         DATOS (Conectados a RMarkdown)                       %
%%%%%%%%%%%%%%%%%%%%%%%%%%%%%%%%%%%%%%%%%%%%%%%%%%%%%%%%%%%%%%%%%%%%%%%%%%%%%%%%
\title{Análisis Descriptivo de los Factores Determinantes en ciclos
agrícolas}
\author{Joshua Santeliz \\ Emily Gallardo \\ }

% Lógica para Facultad: Si la defines en el YAML la usa, si no, usa la default

\facultad{Facultad de Ciencias Económicas y Sociales \\ Escuela de
Estadistica y Ciencias Actuariales}
  \escudofacultad{Latex/Classes/Escudos/faces}

\degree{Computación I}
\director{Jesus Ochoa \\ Oliver Triveño}
\degreedate{Febrero 2026} 
\lugar{Caracas}

\portadatrue 

% Metadatos PDF
\keywords{}
\subject{}

% --- CORRECCIÓN PARA PANDOC NUEVO (Imágenes) ---
\makeatletter
\providecommand{\pandocbounded}[1]{#1}
\makeatother


%%%%%%%%%%%%%%%%%%%%%%%%%%%%%%%%%%%%%%%%%%%%%%%%%%%%%
%                   DOCUMENTO                       %
%%%%%%%%%%%%%%%%%%%%%%%%%%%%%%%%%%%%%%%%%%%%%%%%%%%%%
\begin{document}

\maketitle

%%%%%%%%%%%%%%%%%%%%%%%%%%%%%%%%%%%%%%%%%%%%%%%%%%%%%
%                  PRÓLOGO                          %
%%%%%%%%%%%%%%%%%%%%%%%%%%%%%%%%%%%%%%%%%%%%%%%%%%%%%
\frontmatter

% Puedes agregar dedicatorias desde el YAML si quieres, o dejarlas fijas aquí

%%%%%%%%%%%%%%%%%%%%%%%%%%%%%%%%%%%%%%%%%%%%%%%%%%%%%
%                   ÍNDICES                         %
%%%%%%%%%%%%%%%%%%%%%%%%%%%%%%%%%%%%%%%%%%%%%%%%%%%%%
\setcounter{secnumdepth}{3}
\setcounter{tocdepth}{3}

\tableofcontents



\cleardoublepage

  \cleardoublepage       % Empieza en página derecha (estándar en tesis)
  \chapter*{Resumen}     % Crea el título "Resumen" sin numeración (Capítulo *)
  \addcontentsline{toc}{chapter}{Resumen} % (Opcional) Lo añade al índice
  
  Resumen del trabajo de
investigación\ldots{}             % Aquí se pega el texto que escribiste en el YAML
  
  \cleardoublepage

%%%%%%%%%%%%%%%%%%%%%%%%%%%%%%%%%%%%%%%%%%%%%%%%%%%%%
%                   CONTENIDO (El Corazón)          %
%%%%%%%%%%%%%%%%%%%%%%%%%%%%%%%%%%%%%%%%%%%%%%%%%%%%%
\mainmatter
\def\baselinestretch{1.5}

% --- CAMBIO CRÍTICO 2: Aquí RMarkdown inyecta todo tu texto ---
\chapter*{Introducción}

Introducción del trabajo de investigación

\chapter{Planteamiento del Problema}

En el mundo agrícola hay muchos temas por tratar, desde la salud
nutricional hasta la gastronomía industrial, siendo temas de vital
importancia para la vida humana. En este estudio nos vamos a ubicar en
la agricultura de la India en el ámbito del cultivo, resaltando que una
gestión eficiente de las temporadas de cultivo no solo garantiza
cosechas óptimas, sino que genera seguridad alimentaria y eleva la
calidad de vida de las personas. Analizaremos específicamente las
temporadas de cultivo y como estas son influenciadas por distintos
factores como el clima y los agroquímicos. Estas influencias son de gran
interés para las empresas ya que una descripción de las mismas,
servirían de referencia para mejorar su gestión de inversión.

Por lo antes mencionado surgen problemas por ejemplo a la hora de saber
que cantidad de agroquímicos aplicar o tener de presupuesto para cierta
estación de cultivos y también no saber el nivel de precipitación que
tendremos en un año, lo cual puede ser determinante para prepararse en
los ciclos de cultivo en dicho año. estos problemas se resuelven luego
de modelar esta investigación donde podremos comprender mejor la
situación y se recomendará posibles decisiones a tomar. Teniendo en
cuenta esto surgen las siguientes preguntas de investigación\\

\textbf{1.} ¿De qué manera influyen la cantidad de agroquímicos en los
ciclos de cultivos?\\

\textbf{2.} ¿Por qué son necesario los fertilizantes y los pesticidas en
una temporada de cultivo?\\

\textbf{3.} ¿Cuáles son las precipitaciones anuales que más favorecen a
las temporadas de cultivo?\\

\textbf{4.} ¿La cantidad de agroquímicos varía según los niveles de
precipitación?\\

\section{Justificación}

La justificación de este trabajo de investigación reside en la necesidad
de aplicar técnicas estadísticas rigurosas y herramientas de software
como R y Python para mitigar la incertidumbre de sembrar los cultivos en
el sector agrícola de la India. Desde el punto de vista práctico, los
resultados permitirán a las Industrias tomar decisiones basadas en datos
(Data-Driven Decisions) sobre que materiales invertir o en qué temporada
del año se debe desarrollar. Con una metodología replicable utilizando R
y LaTeX para el análisis exploratorio y descriptivo de grandes volúmenes
de datos.\\

\section{Objetivos}

\subsection{Objetivo General}
\begin{itemize}
  \item{Realizar un análisis estadístico descriptivo de datos de cultivos en diferentes temporadas de la India durante el período 2010-2020.  para identificar y describir las características más influyentes en las temporadas de cultivos.}
\end{itemize}

\subsection{Objetivos Específicos}
\begin{itemize}
  \item{1.      Analizar la variabilidad pluviométrica mediante el uso de indicadores estadísticos, con el fin de determinar una óptima gestión del recurso hídrico.}
  \item{2.      Observar la distribución de los factores que afectan a los cultivos por temporada del año.}
  \item{3.      Categorizar temporadas de mayor y menor demanda tanto de pesticidas como de fertilizantes}
\end{itemize}

\section{Cobertura}

\subsection{Horizontal}

La cobertura horizontal abarca las características de cada cultivo
registrado, donde se consideran las dimensiones descriptivas del cultivo
y las métricas de influencia agrícola.\\

\subsection{Vertical}

La cobertura vertical del trabajo de investigación abarca un total de
\textbf{8405 observaciones} (registros). donde cada fila representa un
cultivo obtenido en la India con sus características correspondientes,
con un alcance temporal que abarca desde el año 2010 hasta el año 2020.
Lo que es suficiente tiempo para dar pie a análisis probabilísticos.

\begin{table}[ht]
\centering
\caption{Variables consideradas en el estudio} 
\resizebox{10cm}{!} {
\begin{tabular}{ccc}
  \hline
\ \ \textbf{TIPO} &  \textbf{VARIABLE} & \textbf{DESCRIPCIÓN} \\ 
  \hline
   Categórica & Crop & Nombre del cultivo \\ 
Categórica &  Season & Temporada que se cultiva \\ 
 Numérica & Crop\_Year & Año que se cultivó (2010-2020) \\
 Numérica & Annual\_Rainfall & La precipitación anual  \\
 Numérica & Fertilizer & Fertilizante utilizado en los cultivos (kg) \\
 Numérica & Pesticide & Pesticida  utilizado en los cultivos (kg)  \\

   \hline
\end{tabular}
}
\end{table}

\section{Periodo de Referencia}

La fuente de datos a estudiar corresponde al repositorio abierto de
Kaggle (``Agricultural Crop Yield in Indian States Dataset''). Donde se
toma como período de referencia los registros pertenecientes al
intervalo de tiempo 2010-2020 , donde la consistencia de los datos es
mayor entre los años 2017-2019, El procesamiento de la información se
realizó con corte a Febrero de 2026.

\section{Variables}

Las variable principal en estudio (Variable Dependiente/Objetivo) es
Season que representan las Temporadas de cultivo, medido en 3 estaciones
``Kharif'' (Otoño), ``Zaid'' (Verano corto) y ''Rabi'' (Invierno)
,siendo las principales temporadas de cultivo en la India. Como
variables explicativas o independientes se analizan:

\begin{enumerate}
  \item Variables Temporales: \textbf{Crop year}
  \item Variables Categóricas: \textbf{Crop} 
  \item Variables influyentes: \textbf{Annual\_Rainfall }, \textbf{Fertilizer } y \textbf{Pesticide }
\end{enumerate}

\chapter{Marco Teórico}

\section{Bases Teóricas}

Resumen de las bases teoricas explicadas en el capitulo y su
importancia\\

\subsection{Análisis Exploratorio de Datos (EDA)}

El Análisis Exploratorio de Datos, formalizado por John Tukey, es el
paso preliminar crítico en cualquier estudio de ciencia de datos. Su
objetivo es resumir las características principales de un conjunto de
datos, a menudo con métodos visuales, sin formular hipótesis \emph{a
priori}.

Matemáticamente, para una variable \(X\) (en este caso,
\texttt{Annual\_Rainfall}), nos interesan sus momentos centrales. Sea
\(x_1, x_2, ..., x_n\) una muestra de precipitaciones observadas,
definimos:\\

\begin{itemize}
\item
  \textbf{Media Aritmética (Media Muestral):} \[
  \bar{x} = \frac{1}{n} \sum_{i=1}^{n} x_i
  \]
\item
  \textbf{Varianza (Volatilidad):} Mide la dispersión de las lluvias
  respecto a la media, fundamental para medir el riesgo de inversión en
  un título. \[S^2 = \frac{1}{n-1} \sum_{i=1}^{n} (x_i - \bar{x})^2\]
\end{itemize}

En el contexto de este trabajo, el EDA permitirá detectar
\emph{outliers} (juegos con ventas excepcionales) y verificar si la
distribución de ventas sigue una normalidad asintótica o presenta colas
pesadas (distribución Pareto/Log-Normal), común en datos financieros.

\subsection{}

Descripción de la definición 2\\

\subsection{Definición n}

Descripción de la definición n\\

\chapter{Marco Metódico}

En esta sección se presentan las metodologías y/o test necesarios
relacionados con la práctica divulgada en el marco teórico, así como las
fuentes de datos.

\section{Bases metódicas utilizadas en el trabajo de investigación}

Desarrollo\\

\subsection{Item 1}

Descripción\\

\subsection{Item 2}

Descripción\\

\subsubsection{Item 2.2}

Descripción\\

\subsection{Item n}

Descripción\\

\chapter{Análisis de Resultados}

\{\small Este capítulo presenta los resultados obtenidos al aplicar las
metodologías descritas en el capítulo anterior, así como del análisis
estadístico descriptivo de las variables estudiadas y las ventajas y
desventajas al aplicar dichos tests.\}

\section{Presentación de resultados 1}

Descripción\\

\section{Presentación de resultados 2}

Descripción\\

\section{Presentación de resultados n}

Descripción\\

\chapter{Conclusiones y Recomendaciones}

\textbf{1.-} Conclusion 1 del trabajo de investigación.\\

\textbf{Recomendación} Recomendación 1 del trabajo de investigación.\\

\textbf{2.-} Conclusion 2 del trabajo de investigación.\\

\textbf{Recomendación} Recomendación 2 del trabajo de investigación.\\

\textbf{n.-} Conclusion n del trabajo de investigación.\\

\textbf{Recomendación} Recomendación n del trabajo de investigación.\\

\appendix
\renewcommand{\thechapter}{\Alph{chapter}}

\chapter{Anexo A}

\chapter{Anexo B}

\chapter{Anexo C}

%%%%%%%%%%%%%%%%%%%%%%%%%%%%%%%%%%%%%%%%%%%%%%%%%%%%%
%                   REFERENCIAS                     %
%%%%%%%%%%%%%%%%%%%%%%%%%%%%%%%%%%%%%%%%%%%%%%%%%%%%%
% Mantenemos la estructura natbib de tu plantilla
\renewcommand{\bibname}{Lista de Referencias}
\bibliographystyle{apalike} 
\bibliography{bibliografia.bib} 

%%%%%%%%%%%%%%%%%%%%%%%%%%%%%%%%%%%%%%%%%%%%%%%%%%%%%
%                   APÉNDICES                       %
%%%%%%%%%%%%%%%%%%%%%%%%%%%%%%%%%%%%%%%%%%%%%%%%%%%%%

\end{document}
